

%% =====================================================================
%%  Figure: One-sided vs Two-sided Capacitated Tree Cut
%% =====================================================================

\begin{figure}[t]
\centering

\begin{subfigure}{0.48\textwidth}
\centering
\scalebox{1.1}{
\begin{tikzpicture}[
  arrow/.style={-{Stealth[length=5pt]}, thick, gray!70},
  lbl/.style={font=\scriptsize\bfseries},
  every node/.style={font=\scriptsize}
]
\def\s{0.65}

%% LEFT: One-sided cut
\begin{scope}[xshift=0cm]

  %% Residual background
  \fill[residual, opacity=0.4, rounded corners=2pt]
    (-0.3*\s,-0.3*\s) rectangle (4*\s+0.3*\s, 3*\s+0.3*\s);

  %% Spanning tree
  \draw[treecol, very thick]
    (0,3*\s) -- (1*\s,3*\s) -- (2*\s,3*\s)
    (3*\s,3*\s) -- (4*\s,3*\s)
    (0,3*\s) -- (0,2*\s) -- (0,1*\s)
    (1*\s,3*\s) -- (1*\s,2*\s) -- (1*\s,1*\s) -- (1*\s,0)
    (2*\s,3*\s) -- (2*\s,2*\s) -- (2*\s,1*\s) -- (2*\s,0)
    (3*\s,3*\s) -- (3*\s,2*\s) -- (3*\s,1*\s) -- (3*\s,0)
    (4*\s,3*\s) -- (4*\s,2*\s) -- (4*\s,1*\s) -- (4*\s,0)
    (0,0) -- (1*\s,0);

  %% Cut edge
  \draw[cutcol, very thick, line cap=round] (2*\s,3*\s) -- (3*\s,3*\s);
  \node[font=\tiny, text=cutcol] at (2.5*\s, 2.3*\s+0.3) {cut};
  \node[font=\tiny, text=cutcol] at (2.1*\s, 2.9*\s+0.3) {$u$};
  \node[font=\tiny, text=cutcol] at (2.9*\s, 3*\s+0.3) {$u^{\prime}$};

  %% Highlight extracted component (right side)
  \fill[distA, opacity=0.3, rounded corners=2pt]
    (3*\s-0.2*\s, -0.2*\s) rectangle (4*\s+0.2*\s, 3*\s+0.2*\s);

  %% Nodes
  \foreach \r in {0,...,3} {
    \foreach \c in {0,...,2} {
      \gnode{\c*\s}{\r*\s}{residual}
    }
    \gnode{3*\s}{\r*\s}{distA}
    \gnode{4*\s}{\r*\s}{distA}
  }

  %% Candidates
  \cand{1*\s}{2*\s}  \cand{4*\s}{1*\s}

  %% Annotations
  \draw[decorate, decoration={brace, amplitude=4pt, mirror}, thick, distA!70!black]
    (4*\s+0.4*\s, -0.2*\s) -- (4*\s+0.4*\s, 3*\s+0.2*\s);
  \node[font=\tiny, distA!70!black, anchor=west] at (4*\s+0.55*\s, 1.5*\s)
    {$T_{u^\prime}$ admissible};

  \draw[decorate, decoration={brace, amplitude=4pt}, thick, gray!60]
    (-0.4*\s, -0.2*\s) -- (-0.4*\s, 3*\s+0.2*\s);
  \node[font=\tiny, gray!60!black, anchor=east, align=center] at (-0.55*\s, 1.5*\s)
    {residual \\continues};

  \node[lbl] at (2*\s, -0.7) {$c^{(r)} > c_{\max}^{\ell}$};
\end{scope}
\end{tikzpicture}
}
\subcaption{One-sided cut}\label{fig:4a}
\end{subfigure}
\hfill
\begin{subfigure}{0.48\textwidth}
\centering
\scalebox{1.1}{
\begin{tikzpicture}[
  arrow/.style={-{Stealth[length=5pt]}, thick, gray!70},
  lbl/.style={font=\scriptsize\bfseries},
  every node/.style={font=\scriptsize}
]
\def\s{0.65}

%% RIGHT: Two-sided cut
\begin{scope}[xshift=0cm]

  %% Component backgrounds
  \fill[distC, opacity=0.3, rounded corners=2pt]
    (-0.2*\s,-0.2*\s) rectangle (2*\s+0.2*\s, 3*\s+0.2*\s);
  \fill[distD, opacity=0.3, rounded corners=2pt]
    (3*\s-0.2*\s,-0.2*\s) rectangle (4*\s+0.2*\s, 3*\s+0.2*\s);

  %% Spanning tree
  \draw[treecol, very thick]
    (0,3*\s) -- (1*\s,3*\s) -- (2*\s,3*\s)
    (3*\s,3*\s) -- (4*\s,3*\s)
    (0,3*\s) -- (0,2*\s) -- (0,1*\s) 
    (1*\s,3*\s) -- (1*\s,2*\s) -- (1*\s,1*\s) -- (1*\s,0)
    (2*\s,3*\s) -- (2*\s,2*\s) -- (2*\s,1*\s) -- (2*\s,0)
    (3*\s,3*\s) -- (3*\s,2*\s) -- (3*\s,1*\s) -- (3*\s,0)
    (4*\s,3*\s) -- (4*\s,2*\s) -- (4*\s,1*\s) -- (4*\s,0)
    (0,0) -- (1*\s,0);

  %% Cut edge
  \draw[cutcol, very thick, line cap=round] (2*\s,3*\s) -- (3*\s,3*\s);
  \node[font=\tiny, text=cutcol] at (2.5*\s, 2.3*\s+0.3) {cut};
  \node[font=\tiny, text=cutcol] at (2.1*\s, 2.9*\s+0.3) {$u$};
  \node[font=\tiny, text=cutcol] at (2.9*\s, 3*\s+0.3) {$u^{\prime}$};

  %% Nodes
  \foreach \r in {0,...,3} {
    \foreach \c in {0,...,2} {
      \gnode{\c*\s}{\r*\s}{distC}
    }
    \gnode{3*\s}{\r*\s}{distD}
    \gnode{4*\s}{\r*\s}{distD}
  }

  %% Candidates
  \cand{1*\s}{2*\s}  \cand{4*\s}{1*\s}

  %% Annotations
  \draw[decorate, decoration={brace, amplitude=4pt, mirror}, thick, distD!70!black]
    (4*\s+0.4*\s, -0.2*\s) -- (4*\s+0.4*\s, 3*\s+0.2*\s);
  \node[font=\tiny, distD!70!black, anchor=west] at (4*\s+0.55*\s, 1.5*\s)
    {$T_{u^{\prime}}$ admissible};

  \draw[decorate, decoration={brace, amplitude=4pt}, thick, distC!70!black]
    (-0.4*\s, -0.2*\s) -- (-0.4*\s, 3*\s+0.2*\s);
  \node[font=\tiny, distC!70!black, anchor=east] at (-0.55*\s, 1.5*\s)
    {$T_u$ admissible};

  \node[lbl] at (2*\s, -0.7) {$c^{(r)} \leq c_{\max}^{\ell}$};

\end{scope}
\end{tikzpicture}
}
\subcaption{Two-sided cut}\label{fig:4b}
\end{subfigure}

\caption{Capacitated tree cut modes. \textbf{(a)}~One-sided mode ($c^{(r)} > c^{\ell}_{\max}$): Feasibility must hold for only side $T_{u^\prime}$, i.e., $\phi(T_{u^\prime})>0$. The residual continues to the next recursion step. \textbf{(b)}~Two-sided mode ($c^{(r)} \leq c^{\ell}_{\max}$): Both sides must satisfy the feasibility condition, i.e., $\phi(T_u)>0$ and $\phi(T_{u^\prime})>0$. Stars (\protect\tikz\protect\node[star, star points=5, star point ratio=2.2, fill=candcol, inner sep=0pt, minimum size=5pt, draw=none] {};) mark candidate facilities.}
\label{fig:cuts}
\end{figure}